\begingroup
\renewcommand{\arraystretch}{1.45}
\fancyhead[L]{\small BESTIMMUNG DER ANZIEHUNG}
\Needspace{41\baselineskip}
\AttractionArticle{6}

Über die Beschaffenheit dieser kubischen Gleichung ist nun Folgendes zu beachten.

I.\quad Aus dem letzten Glied der Gleichung \(-aabbCC\) wird geschlossen, dass sie gewiss eine reelle Wurzel hat, und zwar entweder eine positive oder, wenn \(C=0\), eine der Null gleiche. Bezeichnen wir diese nichtnegative reelle Wurzel mit \(g\).

II.\quad Subtrahiert man von der so dargestellten Gleichung 12
\[
\begin{aligned}
&x=\frac{AAx}{aa+x}+\frac{BBx}{bb+x}+CC
\end{aligned}
\]
diese Gleichung\textperiodcentered
\[
\begin{aligned}
&g=\frac{AAg}{aa+g}+\frac{BBg}{bb+g}+CC
\end{aligned}
\]
und dividiert durch \(x-g\), so entsteht eine neue Gleichung, welche die beiden übrigen Wurzeln enthält
\[
\begin{aligned}
&1=\frac{aaAA}{(aa+x)(aa+g)}+\frac{bbBB}{(bb+x)(bb+g)}
\end{aligned}
\]
und, gehörig geordnet und gelöst, \([14]\) liefert
\[
\begin{aligned}
&2x=\frac{aaAA}{aa+g}+\frac{bbBB}{bb+g}-aa-bb\\[6pt]
&\qquad{}\pm\surd\!\left(\begin{gathered}\left(aa-bb-\frac{aaAA}{aa+g}+\frac{bbBB}{bb+g}\right)^2
\\[6pt]+\frac{4aabbAABB}{(aa+g)(bb+g)}\end{gathered}\right)
\end{aligned}
\]
Da die Größe unter dem Wurzelzeichen ihrer Natur nach positiv oder wenigstens nichtnegativ ist, zeigt dieser Ausdruck, dass auch die beiden übrigen Wurzeln stets reell sind.

\Needspace{28\baselineskip}
III.\quad Subtrahiert man aber die so dargestellten Gleichungen voneinander
\[
\begin{aligned}
&gx=\frac{AAgx}{aa+x}+\frac{BBgx}{bb+x}+gCC\\[6pt]
&gx=\frac{AAgx}{aa+g}+\frac{BBgx}{bb+g}+xCC
\end{aligned}
\]
und dividiert durch \(g-x\), so ergibt sich die Gleichung, welche die beiden übrigen Wurzeln enthält, in dieser Form:
\[
\begin{aligned}
&0=\frac{AAgx}{(aa+g)(aa+x)}+\frac{BBgx}{(bb+g)(bb+x)}+CC
\end{aligned}
\]
die offenbar, wenn \(g\) eine positive Größe ist, durch einen positiven Wert von \(x\) nicht erfüllt werden kann. Daraus schließen wir, dass unsere kubische Gleichung nicht mehr als eine positive Wurzel haben kann.
\SourceLeaf{351}\par
\Needspace{25\baselineskip}
IV.\quad Wann immer also \(0\) nicht unter den Wurzeln unserer Gleichung ist, wird notwendigerweise eine positive Wurzel mit zwei negativen vorhanden sein. Wann immer aber \(C=0\) und somit \(0\) eine der Wurzeln ist, wird die Gleichung
\[
\begin{aligned}
&xx-(AA+BB-aa-bb)x+aabb-aaBB-bbAA=0
\end{aligned}
\]
die übrigen enthalten, woraus diese Wurzeln durch
\[
\begin{aligned}
&\frac12(AA+BB-aa-bb)\pm\frac12\surd\!\left((AA-BB-aa+bb)^2+4AABB\right)
\end{aligned}
\]
ausgedrückt werden. Hier werden wiederum drei Fälle zu unterscheiden sein.

\Needspace{20\baselineskip}
\textit{Erstens}, wenn das letzte Glied \(aabb-aaBB-bbAA\) positiv ist (d.\ h.\ wenn der angezogene Punkt in der Ebene der anziehenden Ellipse \textit{innerhalb} der Kurve liegt), werden beide Wurzeln, da sie reell sein müssen, dasselbe Vorzeichen haben; und da sie nicht beide positiv sein können, werden sie notwendigerweise negativ sein. Dies kann übrigens auch unabhängig von dem schon Bewiesenen daraus geschlossen werden, dass der mittlere Koeffizient, den man so darstellen kann
\[
\begin{aligned}
&(aabb-aaBB-bbAA)\left(\frac1{aa}+\frac1{bb}\right)+\frac{bbAA}{aa}+\frac{aaBB}{bb}
\end{aligned}
\]
in diesem Fall offenbar positiv ist.

\textit{Zweitens}, wenn das letzte Glied negativ ist oder der angezogene Punkt in der Ebene der Ellipse \textit{außerhalb} der Kurve liegt, wird notwendigerweise die eine Wurzel positiv, die andere negativ sein.

\Needspace{18\baselineskip}
\textit{Drittens} aber, wenn das letzte Glied selbst verschwände oder der angezogene Punkt auf dem Umfang der Ellipse selbst läge, würde auch die zweite Wurzel \(=0\) und die dritte
\[
\begin{aligned}
&=-\frac{bbAA}{aa}-\frac{aaBB}{bb}
\end{aligned}
\]
also negativ werden. Diesen physikalisch unmöglichen Fall, in dem die Anziehung selbst unendlich groß würde, werden wir von unserer Untersuchung, wenigstens an dieser Stelle, ausschließen.

\Needspace{46\baselineskip}
\AttractionArticle{7}

Zur Bestimmung der Koeffizienten \(\gamma,\gamma',\gamma''\) finden wir aus den Gleichungen \(1,3,4,6,7,9,10\)
\SourceLeaf{352}
\[
\left.\begin{aligned}
\gamma&=\frac{1}{\surd{\bigl(1-(\frac{aA}{aa+G})^2-(\frac{bB}{bb+G})^2\bigr)}}\\[9pt]
\gamma'&=\frac{1}{\surd{\bigl((\frac{aA}{aa-G'})^2+(\frac{bB}{bb-G'})^2-1\bigr)}}\\[9pt]
\gamma''&=\frac{1}{\surd{\bigl((\frac{aA}{aa-G''})^2+(\frac{bB}{bb-G''})^2-1\bigr)}}
\end{aligned}\right\}\quad[15]
\]
Aus diesen Gleichungen, gehörig mit 5, 8, 11 verbunden, folgt auch:
\[
\left.\begin{aligned}
\gamma&=\surd{\frac{G}{(\frac{AG}{aa+G})^2+(\frac{BG}{bb+G})^2+CC}}\\[9pt]
\gamma'&=\surd{\frac{G'}{(\frac{AG'}{aa-G'})^2+(\frac{BG'}{bb-G'})^2+CC}}\\[9pt]
\gamma''&=\surd{\frac{G''}{(\frac{AG''}{aa-G''})^2+(\frac{BG''}{bb-G''})^2+CC}}
\end{aligned}\right\}\quad[16]
\]
Diese letzteren Ausdrücke zeigen, dass keine der Größen $G,G',G''$ negativ sein kann, wenn $\gamma,\gamma',\gamma''$ reell sein sollen.

\Needspace{34\baselineskip}
In dem Fall also, in dem nicht $C=0$ gilt, muss $G$ notwendigerweise der positiven Wurzel der Gleichung $B$ gleichgesetzt werden; und somit ist offenbar, dass $-G'$ der einen negativen Wurzel und $-G''$ der anderen gleich sein muss\HistoricalNote{26}; welche Wurzel wir aber für $-G'$ und welche für $-G''$ annehmen, wird völlig willkürlich sein.

Wann immer $C=0$ und der angezogene Punkt innerhalb der Kurve liegt, müssen notwendigerweise die beiden negativen Wurzeln der Gleichung 13 für $-G'$ und $-G''$ angenommen und folglich $G=0$ gesetzt werden. Da jedoch in diesem Fall die erste Formel in 16 unbestimmt wird, werden wir an ihrer Stelle die erste Formel in 15 beibehalten, die
\[
\begin{aligned}
&\gamma=\frac{1}{\surd{\bigl(1-\frac{AA}{aa}-\frac{BB}{bb}\bigr)}}
\end{aligned}
\]
liefert. \Needspace{24\baselineskip}
Wann immer aber für $C=0$ der angezogene Punkt außerhalb der Ellipse liegt. ist die positive Wurzel der Glei\SourceLeaf{353}%
chung 13 $=G$ und entweder die negative Wurzel $=-G'$ und $G''=0$ oder die negative Wurzel $=-G''$ und $G'=0$ zu setzen; den Koeffizienten $\gamma''$ oder $\gamma'$ werden wir jedoch durch die Formel
\[
\begin{aligned}
&\frac{1}{\surd{\bigl(\frac{AA}{aa}+\frac{BB}{bb}-1\bigr)}}
\end{aligned}
\]
finden. In dem schon ausgeschlossenen Fall, in dem der angezogene Punkt auf dem Umfang der Ellipse selbst liegend angenommen würde, würden die Koeffizienten $\gamma$ und $\gamma'$ oder $\gamma$ und $\gamma''$ unendlich werden, was anzeigt, dass unsere Transformation auf diesen Fall überhaupt nicht anwendbar ist.

\clearpage
\endgroup
